% TOP_PROBLEM: {"id": 20003177, "title": "A quadratic-relation sieve for formality in dimension seven", "category_id": 7, "category": {"id": 7, "name": "topology", "display_name": "Topology", "description": "Properties preserved under continuous deformations.", "slug": "topology", "order_index": 7, "created_at": "2026-07-31T15:26:25.671Z"}, "status": "partially_solved", "problem_number": "AIM-TOPOLOGY-0265", "set_id": 14, "statusQualification": "Uses rational Sullivan formality, not the stronger and metric-dependent notion of geometric formality."}
% FIRST_POSTED: 2026-10-01
% ENTRY_KIND: statement_only
% UnsolvedMath ID 20003177; source catalogue code AIM-TOPOLOGY-0265
% !TeX program = lualatex
% Definition and sources only; this file is not an LLM solution attempt.
\documentclass[11pt,a4paper]{article}
\usepackage[margin=25mm]{geometry}
\usepackage{fontspec}
\setmainfont{lmroman10-regular.otf}[BoldFont=lmroman10-bold.otf,
 ItalicFont=lmroman10-italic.otf,BoldItalicFont=lmroman10-bolditalic.otf]
\usepackage{amsmath,amssymb,mathtools}
\usepackage{unicode-math}
\setmathfont{latinmodern-math.otf}
\AtBeginDocument{\let\setminus\smallsetminus}
\usepackage{xurl,hyperref}
\hypersetup{colorlinks=true,urlcolor=blue}
\setcounter{secnumdepth}{0}
\setlength{\parindent}{0pt}
\setlength{\parskip}{5pt}
\setlength{\emergencystretch}{3em}
\begin{document}
\section*{A quadratic-relation sieve for formality in dimension seven}\label{20003177}
{\small UnsolvedMath ID 20003177 · AIM-TOPOLOGY-0265 · Topology · AIM Workshop Problem Lists}
\subsection{Definitions and mathematical statement}
Let $M$ be a closed connected smooth manifold admitting a Riemannian metric with strictly positive sectional curvature on every tangent two-plane.
Write $A_{\mathrm{PL}}(M;\mathbb Q)$ for Sullivan's commutative differential graded algebra of rational polynomial differential forms on the singular simplicial set of $M$.
A commutative differential graded algebra has a graded-commutative product and a differential of degree one satisfying the graded Leibniz rule.
Its cohomology inherits a graded-commutative algebra structure.
A morphism is a quasi-isomorphism if it induces an isomorphism in every cohomological degree.

The manifold is rationally formal if there is a finite zigzag of quasi-isomorphisms of such algebras connecting
\[
 A_{\mathrm{PL}}(M;\mathbb Q)
 \quad\text{and}\quad
 (H^*(M;\mathbb Q),0),
\]
where the latter has its ordinary cup product and zero differential.
Equivalently, its Sullivan minimal model is quasi-isomorphic to this zero-differential cohomology algebra.

Does every closed manifold admitting positive sectional curvature have this rational formality property?
The conclusion concerns the algebraic homotopy information of the manifold and is independent of the chosen positively curved metric.
It is distinct from geometric formality, which requires wedge products of harmonic forms for a particular metric to remain harmonic.
No such harmonic-form condition is requested here.
The question permits arbitrary dimension and imposes neither homogeneity nor a prescribed symmetry action.
A negative answer requires a positively curved manifold whose rational differential graded algebra cannot be reduced, up to quasi-isomorphism, to its cohomology ring alone.
\subsection{Short English statement}
Does positive sectional curvature force Sullivan rational formality of the underlying closed manifold?
\subsection{Sources}
\begin{itemize}
\item[S1] ULAM.AI and the UnsolvedMath Contributors, supplied problems.json, record 20003177 (AIM-TOPOLOGY-0265), AIM Workshop Problem Lists. Catalogue identity and source transcription; CC BY 4.0.
\url{https://huggingface.co/datasets/ulamai/UnsolvedMath}.
\item[S2] American Institute of Mathematics, Manifolds with nonnegative sectional curvature, item 52. Original workshop question underlying AIM-TOPOLOGY-0265.
\url{https://aimath.org/WWN/nnsectcurvature/nnsectcurvature.pdf}.
\end{itemize}
\end{document}
